\section{The double Fock space of type B}
%\label{subsec:Fock}

Let $H_\R$ be a separable real Hilbert space and let $H$ be its
complexification with the inner product~$\langle\cdot,\cdot\rangle$, linear in the right component and anti-linear in the left.
The Hilbert space $\HH:=H\otimes {H}$ is the complexification of its
real subspace $\HH_\R:=H_\R\otimes {H}_\R$, with the inner product
\[
\langle x\otimes y,\xi\otimes \eta \rangle_{\HH} = \langle
x,\xi\rangle\langle y,\eta\rangle.
\]

We define $\H:=H^{\otimes n} \otimes {{H}^{\otimes n}}=H^{\otimes 2n}$ and instead of
indexing its simple tensors by $\{1,\dots,2n\}$ we will index them by $[\pm n]=\{-n,\dots,-1,1,\dots,n\}$ and we will use the
typical involution $\bar n=-n$ for $n\in \N$
\[
\H\ni\mathbf{x}_{\overline n} \otimes \mathbf{x}_n=x_{\overline n}\otimes\cdots \otimes x_{\overline 1} \otimes x_1\otimes\cdots \otimes x_n=x_{\overline n}\otimes \cdots \otimes x_{n}.
\]
We use this convention for indexing the elements of $\H$ to define a natural action of the
hyperoctahedral group $B(n)$ on $\H$ by setting
\begin{align*}
\sigma\colon\, &\H\to \H,
\\
	&x_{\overline n}\otimes\cdots \otimes x_{\overline 1} \otimes x_1\otimes\cdots \otimes x_n\mapsto x_{\sigma^{-1}(\overline n)}\otimes\cdots \otimes x_{\sigma^{-1}(\overline 1)} \otimes x_{\sigma^{-1}(1)}\otimes\cdots \otimes x_{\sigma^{-1}(n)}
\end{align*}
for any $\sigma \in B(n)$.

\begin{Remark}
	It is worth to mention that in the present paper we act on the indexes, but in the previous one, \cite{BozejkoEjsmontHasebe2015}, we were acting on the vectors.
\end{Remark}



Let $\F$ be the algebraic full Fock space over $\HH$:
\begin{equation*}
	\F:= \bigoplus_{n=0}^\infty \H= \bigoplus_{n=0}^\infty H^{\otimes 2n},
\end{equation*}
with the convention that $\HH^{\otimes 0} =H^{\otimes 0} \otimes H^{\otimes 0}=\C\Omega \otimes \Omega$ is the one-dimensional normed space along with the unit vector $\Omega \otimes \Omega$.
We equip $\F$ with the inner product
\[
\langle x_{\overline n} \otimes \cdots \otimes x_{\overline 1} \otimes x_1 \otimes \cdots \otimes x_n, y_{\overline m} \otimes \cdots \otimes y_{\overline 1} \otimes y_1 \otimes \cdots \otimes y_m\rangle_{0,0}:= \delta_{m,n}\prod_{i=\overline n }^n \langle x_i, y_i\rangle.
\]
By definition, the action of the generators $\pi_i$ for $0\leq i \leq n-1$ on
\begin{align*}
\eta = x_{\bar n} \otimes \cdots \otimes x_{\bar 1} \otimes x_1 \otimes \cdots \otimes x_{n}\in \H
\end{align*}
is given for $i\geq 1$ and $n\geq 2$ by
\begin{align*}
\pi_i(\eta) =
	x_{\bar n} \otimes \cdots \otimes x_{\overline{i}} \otimes x_{\overline{i+1}} \otimes \cdots \otimes x_{\bar 1} \otimes x_1 \otimes \cdots \otimes x_{i+1} \otimes x_{i} \otimes \cdots \otimes x_{n}, \end{align*}
and for $i=0$ and $n\geq 1$ by
\begin{align*}
\pi_0(\eta) = x_{\bar n} \otimes \cdots \otimes x_1 \otimes x_{\bar 1}\otimes \cdots \otimes x_{n}.
\end{align*}

For the parameters $ -1\leq \alpha, q \leq 1$, we define the symmetrization operators
\begin{align*}
	&P_{\s,\q}^{(n)}:= \sum_{\sigma \in B(n)}\s^{l_1(\sigma)} \q^{l_2(\sigma)} \, \sigma,\qquad n \geq1, \\
	&P_{\s,\q}^{(0)}:= \id_{H^{\otimes 0}\otimes H^{\otimes 0}}.
	\intertext{Moreover, let }
	&P_{\s,\q}:=\bigoplus_{n=0}^\infty P_{\s,\q}^{(n)}
\end{align*} be the \emph{type B symmetrization operator} acting on the algebraic full Fock space.
\begin{Remark}
	From
	Bo\.zejko and Speicher \cite[Theorem 2.1]{BozejkoSpeicher1994}, the operator $P_{\s,\q}^{(n)}$ is positive for $ -1\leq \alpha, q \leq 1$. If $ -1< \alpha, q <1$, then $P_{\s,\q}^{(n)}$ is a strictly positive operator meaning that it is positive and
	$\ker P_{\s,\q}^{(n)}=\{0\}$.
\end{Remark}
For $\x\in \HH_{ n}$ and $\mathbf{y}_{\overline m} \otimes \mathbf{y}_m\in
\HH_{ m}$, we deform the inner product $\langle\cdot,\cdot\rangle_{0,0}$ by using the type~B symmetrization operator:
\begin{align*}
	\langle \x,\mathbf{y}_{\overline m} \otimes \mathbf{y}_m\rangle_{\s,\q}:=\delta_{n,m}\langle \x,P_{\s,\q}^{(m)}\mathbf{y}_{\overline m} \otimes \mathbf{y}_m\rangle_{0,0}.
\end{align*}

For $x\in H$, let $l(x)$ and $r(x)$ be the free left and right
annihilation operators on $H^{\otimes n} $, respectively, defined by
the equations
\begin{gather*}
l^\ast(x)(x_1\otimes \dots \otimes x_n ):=x\otimes x_1\otimes \dots \otimes x_n ,
\\
l(x)(x_1\otimes \dots \otimes x_n ):=\langle x, x_1 \rangle x_2\otimes \dots \otimes x_n ,
\\
r^\ast(x)(x_1\otimes \dots \otimes x_n ):= x_1\otimes \dots \otimes x_n \otimes x,
\\
r(x)(x_1\otimes \dots \otimes x_n ):=\langle x, x_n \rangle x_1\otimes \dots \otimes x_{n-1},
\end{gather*}
where the adjoint is taken with respect to the free inner product.
The left-right creation and annihilation operators $\r^\ast(x\otimes y)$, $\r(x\otimes y)$ on $\F$ are defined as follows:
\begin{alignat*}{2}
&\r^\ast(x\otimes y)(\x):=l^\ast(x)\mathbf{x}_{\overline n}\otimes
	r^\ast( y) \mathbf{x}_n, \qquad &&\r^\ast(x\otimes y)\Omega \otimes \Omega:=x \otimes y,
\\
&\r(x\otimes y)(\x ):= l( x) \mathbf{x}_{\overline n}\otimes r(y)\mathbf{x}_n , \qquad &&\r(x\otimes y)\Omega \otimes \Omega :=0,
\end{alignat*}
where $\x\in \H$, $n\geq 1$.
It holds that $[\r^\ast(x\otimes y)]^\ast = \r(x\otimes y)$ where the adjoint is taken with respect to $\langle \cdot, \cdot \rangle_{0,0}$.

\begin{Definition} Let $\q,\s \in (-1,1)$. The algebraic full Fock space $\F$ equipped with the inner product $\langle\cdot,\cdot \rangle_{\s,\q}$ is called the \emph{double Fock space of type B} and denoted by $\mathcal{F}_{\s,\q}(\HH)$.
	For $x\otimes y \in \HH$ we define $\B^\ast(x\otimes y):=
	\r^\ast(x\otimes y)$ and we consider its adjoint operator $\B(x\otimes
	y)$ with
	respect to the inner product $\langle\cdot,\cdot \rangle_{\s,\q}$
	acting on the Hilbert space $\mathcal{F}_{\s,\q}(\HH)$. The operators $\B^\ast(x\otimes y)$ and $\B(x\otimes y)$ are called \emph{double creation and double annihilation operator of type B, respectively}.
\end{Definition}

\begin{Remark}
We would like to emphasize that our double Fock space of type B is different from~\cite{BozejkoEjsmontHasebe2015}.
	In the present version of the model, we apply a more natural operation on tensor product, which finally contributes to the more natural combinatorics of calculating the Gaussian moments.
	In~\cite{BozejkoEjsmontHasebe2015}, we put the action of the symmetrization on $n$ points but in the current ap\-proach we give $2n$ points which is close to our construction from the articles~\cite{BDEG2021,Ejsmont2020}. In~\cite{BozejkoEjsmontHasebe2015}, we defined the creator in the same way as in the case of type A, and then we did not obtain a~natural combinatorics which is compatible with partitions of the set $\{\pm 1,\dots, \pm n\}$.
	Now if we start our construction with the creator operator of double type B, the situation is completely different from~\cite{BozejkoEjsmontHasebe2015}.
\end{Remark}
The following proposition can be derived directly
from \cite[Proposition 2.3]{BozejkoEjsmontHasebe2015}.

\begin{Proposition}%\label{prop1}
We have the decomposition
\begin{equation*}%\label{decomposition}
P^{(n)}_{\s,\q}=\big( \id \otimes P^{(n-1)}_{\s,\q}\otimes \id \big)R^{(n)}_{\s,\q}\qquad \text{on}\quad \H, \qquad n\geq 1,
\end{equation*}
where\vspace{-2mm}
\begin{gather*}
R^{(n)}_{\s,\q} = \id+\sum_{k=1}^{n-1}\q^{k}\pi_{n-1}\cdots \pi_{n-k} + \s \q^{n-1}\pi_{n-1} \pi_{n-2} \cdots \pi_{1}\pi_0\bigg(1+\sum_{k=1}^{n-1}\q^{k}\pi_{1}\cdots \pi_{k}\bigg), \quad n\geq 2,
\\
R^{(1)}_{\s,\q} =\id+\s\pi_0.% \quad n=1.
\end{gather*}
\end{Proposition}

We can compute the annihilation operator in terms of $R_{\s,\q}^{(n)}$ by using the similar
method to that in \cite[Proposition~2.4]{BozejkoEjsmontHasebe2015}.

\begin{Proposition}%\label{prop2}
For $n \geq1$ and $x\otimes y\in \HH_\R$, we have\vspace{-1mm}
\begin{equation*}
\B(x\otimes y)=\r(x\otimes y) R^{(n)}_{\s,\q}\qquad \text{on}\quad \H.
\end{equation*}
\end{Proposition}

\begin{Remark}\label{thm1}
	Using the above notation, we can decompose $\B(x\otimes y)$ into the positive part $p_\q$ and the negative part $\ell_{\q}$ as\vspace{-1mm}
\[
	\B(x\otimes y)= p_\q(x\otimes y)+ \s \ell_{\q}(x\otimes y), \qquad x\otimes y \in \HH,
\]
where\vspace{-2mm}
\begin{align}\label{rq}
&p_\q(x\otimes y)\eta = \sum_{k=1}^n \q^{n-k}\langle x, x_{\bar{k}}\rangle \langle y, x_k \rangle\, x_{\bar n }\otimes \cdots \otimes \check{x}_{\bar k} \otimes \cdots \otimes x_{\bar 1} \nonumber
\\[-1mm]
&\hphantom{p_\q(x\otimes y)\eta = \sum_{k=1}^n}
{}\otimes x_1\otimes \cdots \otimes \check{x}_k \otimes \cdots \otimes x_n,
\\
&\ell_{\q}(x\otimes y)\eta =\q^{n-1}\sum_{k=1}^n \q^{k-1}\langle x, x_{{k}}\rangle \langle y, x_{\bar{k}}\rangle\, x_{\bar n }\otimes \cdots \otimes \check{x}_{\bar k} \otimes \cdots \otimes x_{\bar 1} \nonumber
\\
&\hphantom{\ell_{\q}(x\otimes y)\eta =\q^{n-1}\sum_{k=1}^n}
\otimes x_1\otimes \cdots \otimes \check{x}_k \otimes \cdots \otimes x_n,\label{lq}
	\end{align}
	where $\eta=x_{\bar n }\otimes \cdots \otimes x_n$.
	We called the operator $p_\q$ the positive part and $\ell_{\q}$ the negative part because in the next section we see that they contribute to positive and negative partitions of type B, respectively.
\end{Remark}
